\section{Prospective Comparative Human Evaluation}
\label{app:human}
The revised, separately packaged questionnaire compares written world designs
under blinded source labels. One arm pairs successful outputs of different
recorded models under the same one-sentence prompt (17 artifacts, 41 pairs).
A second arm pairs compositional Agora with single-pass outputs under ten
identical archived requests, using the first replicate of each method.
Source fields, ordering, item quotas, and truncation rules are fixed across
conditions; no summary is rewritten to favor one output. Unavailable model
artifacts are recorded as generation failures, not imputed preference losses.

Each participant reviews two pairs per arm. Presentation order and A/B
orientation are balanced in an 82-code schedule; this is a coverage plan,
not a power calculation. The primary question asks overall design preference.
Five separate diagnostic items address premise realization, rules and
activities, social connections, places and objects, and possible actions.
Each permits a directional preference, a tie, or insufficient evidence.
Comments and quotation permission are optional. Instructions are bilingual,
while source evidence remains in its original English, restricting recruitment
to readers of English. No cross-language equivalence claim is made.

The prospective analysis reports prompt-specific outcomes, ties, abstentions,
and participant-clustered intervals conditional on the fixed corpus. It does
not treat repeated ratings as independent participants or aggregate different
constructs into a validated scale. The unequal-compute architecture comparison
estimates preferences for the observed system configurations. No matching
baseline images exist, so image quality and actual play experience are outside
the comparative instrument. Clear construct and reporting boundaries follow
human-evaluation guidance \cite{howcroft}. No participant count, preference,
agreement estimate, or completed human-study result is claimed in this paper.
